\chapter{Robustness to Global Perturbations}
\label{ch:rq2}

This chapter investigates how global Gaussian noise affects reconstructions and
token assignments. We assess robustness through reconstruction quality and
changes in token assignments, then aggregate assignments across the dataset to
measure changes in code frequencies. The autoregressive model receives token
sequences, so a reconstruction may remain recognizable even when many
assignments change. Nearest-neighbor quantization can switch an assignment when
an encoder output crosses a codebook boundary, although this experiment does
not measure distances to those boundaries.

\section{Experimental Setup}
\label{sec:rq2_setup}

\subsection{Perturbation Protocol}
\label{subsec:rq2_perturbation_protocol}

Let $x \in [-1,1]^{H \times W \times 3}$ denote an input image in the
tokenizer's internal range. For a noise level $\sigma$, the perturbed image is
defined by:
\begin{equation}
    x_{\sigma} = \mathrm{clip}(x + \sigma \epsilon, -1, 1),
    \qquad
    \epsilon \sim \mathcal{N}(0, I).
    \label{eq:rq2_global_noise}
\end{equation}
Clipping keeps the perturbed image in the allowed range and can place values at
its limits. The experiment uses $\sigma\in\{0.1,0.2,0.5\}$ together with the
clean baseline $\sigma=0$. Images are converted to $[0,1]$ for reported image
metrics.

For a tokenizer with encoder $E$, quantizer $q$, and decoder $D$, the perturbed
reconstruction and tokenization are given by:
\begin{equation}
    \hat{x}_{\sigma} = D(q(E(x_{\sigma}))),
    \qquad
    z_{\sigma} = q(E(x_{\sigma})).
    \label{eq:rq2_reconstruction_tokenization}
\end{equation}
Applying the same operations to $x$ gives the clean reconstruction
$\hat{x}_0=D(q(E(x)))$ and clean token grid $z_0=q(E(x))$. Each $z_{\sigma}$ is
a $16\times16$ grid.

The experiment was performed on the full ImageNet validation set for VQGAN and
LlamaGen. Each image is evaluated clean and at $\sigma=0.1$, $0.2$, and $0.5$.

Both tokenizers use the same images and nominal noise levels, but their runs use
different Gaussian noise realizations. Model comparisons are therefore matched
by image and noise level, not by individual noise samples.

\subsection{Measured Effects}
\label{subsec:rq2_metrics}

\paragraph{Reconstruction fidelity.}
PSNR measures the input corruption, $\mathrm{PSNR}(x,x_{\sigma})$; fidelity to
the noisy input, $\mathrm{PSNR}(x_{\sigma},\hat{x}_{\sigma})$; and fidelity to
the clean input, $\mathrm{PSNR}(x,\hat{x}_{\sigma})$. The last quantity is the
main reconstruction measure in this chapter.

\paragraph{Token flip rate.}
The token flip rate is the fraction of assignments that differ from the clean
tokenization:
\begin{equation}
    \mathrm{flip}_{\sigma}(x)
    =
    \frac{1}{N}
    \sum_{i=1}^{N}
    \mathbf{1}\!\left[z_{\sigma,i} \neq z_{0,i}\right],
    \label{eq:rq2_token_flip}
\end{equation}
where $N=256$ is the number of positions in the $16 \times 16$ token grid. A value near zero indicates that most positions keep the same token under noise; a value near one indicates that nearly all token identities changed.

Flip rates are directly comparable across noise levels within a tokenizer.
Their absolute values are only descriptive across tokenizers: LlamaGen uses
almost the full nominal vocabulary, whereas VQGAN quantizes into a much
smaller empirical subset, and the two codebooks induce different partitions
of latent space. A higher cross-model flip rate therefore does not by itself
imply a worse representation.

\paragraph{Codebook usage.}
After measuring changes at individual positions, we aggregate assignments
across images to examine code-frequency changes. For one empirical distribution
$p_k$, the entropy and corresponding perplexity are given by:
\begin{equation}
    H(p) = -\sum_k p_k \log p_k,
    \qquad
    \mathrm{PPL}(p) = \exp(H(p)).
    \label{eq:rq2_perplexity}
\end{equation}
Active count measures codebook coverage. Top-$k$ mass measures how strongly
assignments concentrate into a small number of codes.

\section{Experimental Results}
\label{sec:rq2_results}

\subsection{Qualitative Behavior}
\label{subsec:rq2_qualitative}

Figure~\ref{fig:rq2_qual_combined} shows selected
ImageNet examples for both tokenizers. Each example contains the clean image,
inputs at $\sigma=0.1$, $0.2$, and $0.5$, the clean reconstruction, and the
three perturbed reconstructions. These panels are qualitative examples; the
quantitative conclusions use the full validation set.

\begin{figure}[p]
    \centering
    \textbf{(a) VQGAN}\\[2pt]
    \includegraphics[width=0.94\textwidth,height=0.36\textheight,keepaspectratio]{ch05_fig_global_noise_qualitative_vqgan.png}\\[6pt]
    \textbf{(b) LlamaGen}\\[2pt]
    \includegraphics[width=0.94\textwidth,height=0.36\textheight,keepaspectratio]{ch05_fig_global_noise_qualitative_llamagen.png}
    \caption{Selected reconstructions under global Gaussian noise. Both panels
    use the same ImageNet samples and show $\sigma=0.1$, $0.2$, and $0.5$.}
    \label{fig:rq2_qual_combined}
\end{figure}
\subsection{Reconstruction Degradation}
\label{subsec:rq2_image_space}

Table~\ref{tab:rq2_psnr_summary} reports $\mathrm{PSNR}(x,\hat{x}_{\sigma})$, the end-to-end fidelity of the noisy-input reconstruction to the original clean image. Figure~\ref{fig:rq2_psnr_x_xhat} visualizes the same trend. Both tokenizers degrade monotonically as $\sigma$ increases, and LlamaGen remains above VQGAN at every level. At the clean baseline, LlamaGen reaches $20.79$ dB and VQGAN reaches $19.65$ dB. At the highest perturbation level, LlamaGen remains at $16.22$ dB, while VQGAN falls to $14.90$ dB.

The clean reconstruction difference persists under the tested perturbations.
The gaps are 1.14, 1.09, 1.00, and 1.32 dB from $\sigma=0$ to $0.5$, so they are
broadly similar rather than monotonically widening.

\begin{table}[t]
    \centering
    \small
    \begin{tabular}{llrr}
        \textbf{Tokenizer} & \textbf{Noise level} & $\boldsymbol{\sigma}$ & $\boldsymbol{\mathrm{PSNR}(x,\hat{x}_{\sigma})}$ \\ \hline
        LlamaGen & clean & 0.0 & $20.79 \pm 3.18$ \\
        LlamaGen & low & 0.1 & $20.32 \pm 2.71$ \\
        LlamaGen & mid & 0.2 & $19.11 \pm 1.98$ \\
        LlamaGen & high & 0.5 & $16.22 \pm 0.94$ \\
        \hline
        VQGAN & clean & 0.0 & $19.65 \pm 3.41$ \\
        VQGAN & low & 0.1 & $19.23 \pm 2.92$ \\
        VQGAN & mid & 0.2 & $18.11 \pm 2.27$ \\
        VQGAN & high & 0.5 & $14.90 \pm 1.03$ \\
    \end{tabular}
    \caption{End-to-end reconstruction fidelity under global Gaussian perturbations. Values are mean PSNR in dB $\pm$ standard deviation over 50{,}000 ImageNet validation images.}
    \label{tab:rq2_psnr_summary}
\end{table}
\begin{figure}[t]
    \centering
    \includegraphics[width=0.86\textwidth]{ch05_fig_psnr_x_xhat_vs_sigma.png}
    \caption{Reconstruction fidelity to the clean input. Points show mean
    $\mathrm{PSNR}(x,\hat{x}_{\sigma})$ and error bars show one standard
    deviation across images. The horizontal axis is $\sigma$ and PSNR is in dB.}
    \label{fig:rq2_psnr_x_xhat}
\end{figure}
Figure~\ref{fig:rq2_psnr_xnoisy_xhat} reports
$\mathrm{PSNR}(x_{\sigma},\hat{x}_{\sigma})$. At $\sigma=0.1$, the values are
$19.26$ dB for LlamaGen and $18.36$ dB for VQGAN, a $0.90$ dB gap. At
$\sigma=0.5$, they are $11.87$ and $11.31$ dB, a $0.56$ dB gap. LlamaGen has a
modest advantage under this measure.

\begin{figure}[t]
    \centering
    \includegraphics[width=0.86\textwidth]{ch05_fig_psnr_xnoisy_xhat_vs_sigma.png}
    \caption{PSNR between noisy input $x_{\sigma}$ and reconstruction $\hat{x}_{\sigma}$. This separates fidelity to the corrupted input from fidelity to the original clean image.}
    \label{fig:rq2_psnr_xnoisy_xhat}
\end{figure}
\subsection{Token Instability}
\label{subsec:rq2_token_flip}

The degradation in image space is gradual, while quantization can cause abrupt
assignment changes. Table~\ref{tab:rq2_token_flip} and
Figure~\ref{fig:rq2_token_flip} report the dataset mean of
$\mathrm{flip}_{\sigma}(x)$. At $\sigma=0.1$, VQGAN changes $75.1\%$ of token
positions and LlamaGen changes $79.5\%$, even though the selected
reconstructions remain recognizable. The latter is a qualitative observation.
At $\sigma=0.2$, the mean flip rates rise to $89.0\%$ and $91.0\%$; at
$\sigma=0.5$, both exceed $97\%$.

LlamaGen reconstructs better in image space, yet its mean flip rate is slightly
higher than VQGAN's at every tested level. This implies that reconstruction
quality and assignment stability describe different responses to noise.

\begin{table}[t]
    \centering
    \small
    \begin{tabular}{llrr}
        \textbf{Tokenizer} & \textbf{Noise level} & $\boldsymbol{\sigma}$ & \textbf{Token flip rate} \\ \hline
        LlamaGen & low & 0.1 & $0.795 \pm 0.099$ \\
        LlamaGen & mid & 0.2 & $0.910 \pm 0.054$ \\
        LlamaGen & high & 0.5 & $0.979 \pm 0.017$ \\
        VQGAN & low & 0.1 & $0.751 \pm 0.103$ \\
        VQGAN & mid & 0.2 & $0.890 \pm 0.057$ \\
        VQGAN & high & 0.5 & $0.975 \pm 0.017$ \\
    \end{tabular}
    \caption{Fraction of token positions changed relative to the clean tokenization. Values are mean $\pm$ standard deviation over ImageNet validation images.}
    \label{tab:rq2_token_flip}
\end{table}
\begin{figure}[t]
    \centering
    \includegraphics[width=0.86\textwidth]{ch05_fig_token_flip_vs_sigma.png}
    \caption{Token flip rate under global Gaussian perturbations. The discrete token representation changes sharply even when the image-space reconstruction remains visually structured.}
    \label{fig:rq2_token_flip}
\end{figure}
\subsection{Global Codebook Distribution}
\label{subsec:rq2_codebook_distribution}

After measuring changes at individual positions, we aggregate assignments
across images to examine code-frequency changes.
Table~\ref{tab:rq2_codebook_distribution} summarizes active-code counts,
perplexity, and top-$k$ concentration; Figure~\ref{fig:rq2_perplexity}
visualizes the main trend. Figure~\ref{fig:rq2_code_frequency_comparison}
compares clean and $\sigma=0.5$ assignment frequencies for the same codes.

The two tokenizers have different baselines. LlamaGen uses the full codebook on
clean ImageNet, with 16{,}384 active codes and a perplexity of 16{,}179. VQGAN
uses 972 active codes, with a perplexity of 910. Under noise, LlamaGen's active
count remains high and reaches 16{,}203 at $\sigma=0.5$; VQGAN's stays near 972.

Perplexity is more sensitive than active-code count. LlamaGen's perplexity decreases from 16{,}179 to 6{,}434 as $\sigma$ increases from 0 to 0.5. VQGAN's perplexity decreases from 910 to 560. Thus, both distributions become more concentrated even when the number of active codes changes little. The top-100 mass confirms this concentration: LlamaGen increases from 0.012 to 0.115, while VQGAN increases from 0.165 to 0.388.

\begin{table}[t]
    \centering
    \small
    \begin{tabular}{llrrrr}
        \textbf{Tokenizer} & \textbf{Level} & \textbf{Active codes} & \textbf{Perplexity} & \textbf{Top-100 mass} & \textbf{Top-1000 mass} \\ \hline
        LlamaGen & clean & 16{,}384 & 16{,}179 & 0.012 & 0.084 \\
        LlamaGen & low & 16{,}381 & 12{,}829 & 0.040 & 0.198 \\
        LlamaGen & mid & 16{,}369 & 10{,}388 & 0.063 & 0.269 \\
        LlamaGen & high & 16{,}203 & 6{,}434 & 0.115 & 0.416 \\
        \hline
        VQGAN & clean & 972 & 910 & 0.165 & 1.000 \\
        VQGAN & low & 973 & 859 & 0.202 & 1.000 \\
        VQGAN & mid & 971 & 780 & 0.249 & 1.000 \\
        VQGAN & high & 972 & 560 & 0.388 & 1.000 \\
    \end{tabular}
    \caption{Global codebook distribution statistics under global perturbations. Top-$k$ mass is the fraction of all token assignments explained by the $k$ most frequent codes at that noise level.}
    \label{tab:rq2_codebook_distribution}
\end{table}
\begin{figure}[t]
    \centering
    \includegraphics[width=0.86\textwidth]{ch05_fig_perplexity_vs_sigma.png}
    \caption{Effective codebook perplexity under global perturbations. LlamaGen remains much broader than VQGAN, but its effective vocabulary contracts substantially as noise increases.}
    \label{fig:rq2_perplexity}
\end{figure}
\begin{figure}[t]
    \centering
    \includegraphics[width=0.96\textwidth]{ch05_fig_code_frequency_clean_vs_sigma_0_5.png}
    \caption{Per-code assignment frequencies for clean inputs and $\sigma=0.5$. Within each tokenizer, codes are ordered once by their clean frequency; the noisy points retain that fixed code ordering.}
    \label{fig:rq2_code_frequency_comparison}
\end{figure}
\section{Discussion}
\label{sec:rq2_discussion}

\paragraph{Reconstruction and assignment response.}
\label{subsec:rq2_reconstruction_token_divergence}

The two tokenizers show different responses in reconstruction quality and token
assignments. LlamaGen has higher PSNR at every tested level, yet its token flip
rate is slightly above VQGAN's.

One possible explanation is codebook geometry, but this experiment does not
measure feature displacements or distances to codebook boundaries. Measuring
those quantities is future work.

These assignment changes alter the sequences supplied to an autoregressive
model. Their effect on generator robustness remains a hypothesis. The
observations suggest that exact assignments can be sensitive even when selected
reconstructions remain recognizable; they do not establish whether individual
tokens carry stable semantics.

\paragraph{Code-frequency concentration.}
\label{subsec:rq2_noise_concentrates_usage}

Perplexity decreases and top-100 mass increases for both tokenizers while
active-code counts change comparatively little. The frequency distributions
therefore become more concentrated under noise.

The results above provide the reference point for the locality analyses. Global
noise perturbs all image regions. Chapter~\ref{ch:rq3} isolates a local image
change, and Chapter~\ref{ch:rq4} studies bounded token edits. These experiments
separate where a response occurs, but do not identify individual causal
mechanisms.
